\documentclass[11pt,a4paper]{article}
\usepackage[utf8]{inputenc}
\usepackage[margin=1in]{geometry}
\usepackage{booktabs}
\usepackage{longtable}
\usepackage{hyperref}
\usepackage{xcolor}
\usepackage{amsmath}
\usepackage{parskip}
\usepackage{enumitem}
\hypersetup{colorlinks=true,urlcolor=blue,linkcolor=black}

\title{Replication Report: Pal, Sharma \& Subramanian (2021)\\
\large Complete genome sequence and identification of polyunsaturated fatty acid biosynthesis genes of the myxobacterium \textit{Minicystis rosea} DSM 24000\textsuperscript{T}}
\author{Ollie (OpenClaw subagent) \\ Set: BVBRC-51}
\date{Replication date: 2026-07-02}

\begin{document}
\maketitle

\begin{center}
\fbox{\parbox{0.85\textwidth}{\centering
\textbf{Paper:} Pal S, Sharma G, Subramanian S. \textit{BMC Genomics} 22:655 (2021). \\
\textbf{DOI:} 10.1186/s12864-021-07955-x \quad \textbf{PMID:} 34511070 \quad \textbf{PMC:} PMC8436480 \\
\textbf{Open access:} CC BY 4.0 (BMC) \\[4pt]
\textbf{Compute:} CherryRd (BLAST+, Python) + uicgpu (antiSMASH 8.0.4). Free endpoints only. \\
\textbf{Verdict:} \textbf{PARTIAL REPLICATION (strong)}
}}
\end{center}

\section{Paper}

A complete-genome announcement for the soil myxobacterium \textit{Minicystis rosea} DSM 24000\textsuperscript{T} (suborder Sorangiineae, family Polyangiaceae). Using PacBio, the authors assembled a \textbf{single circular 16.04\,Mbp chromosome} --- reported as the \textbf{largest bacterial genome sequenced to date (2021)} --- with $\sim$44\% paralogous coding potential. They mine the genome with \textbf{antiSMASH v5.0} (47 biosynthetic gene clusters, BGCs) and specifically identify and phylogenetically place the \textbf{polyunsaturated fatty acid (\textit{pfa}) biosynthetic gene cluster}, arguing for its acquisition via horizontal gene transfer from Actinobacteria based on conserved \textit{pfa} synteny.

This maps cleanly to a BV-BRC Comprehensive Genome Analysis workflow (assembly stats + RASTtk-style annotation) plus secondary-metabolite / specialty-gene detection (antiSMASH).

\section{Claims Tested}

\begin{longtable}{@{}c p{9cm} l c@{}}
\toprule
\# & Claim & Type & Tested? \\
\midrule
C1 & Genome deposited publicly (CP016211.1 / PRJNA321464). & Data availability & \checkmark \\
C2 & Single complete circular chromosome, 16{,}040{,}666\,bp, 69.07\% GC; largest bacterial genome (2021). & Genome stats & \checkmark \\
C3 & 14{,}018 CDS (6{,}983 +, 7{,}035 $-$), 88 tRNA, 4 rRNA operons, coding density 87.31\%. & Annotation & \checkmark \\
C4 & antiSMASH $\to$ 47 BGCs; dominant NRPS/terpene/PKS/RiPP + rare singletons. & Secondary metabolism & \checkmark \\
C5 & \textit{pfa} PUFA cluster present; Pfa1/2/3 homologous to \textit{Aetherobacter}/\textit{Sorangium} Pfa, conserved synteny. & Comparative genomics & \checkmark \\
\bottomrule
\end{longtable}

\section{Method}

All data are free/public; all tools are open-source. Detailed URLs, accessions, and checksums are recorded in \texttt{artifact\_harvest.md}; the chronological record is in \texttt{attempt\_log.md}.

\begin{enumerate}[leftmargin=1.2cm]
\item \textbf{Paper harvest.} Europe PMC full-text XML for PMC8436480; extracted accession CP016211.1 and BioProject PRJNA321464 and the antiSMASH/PUFA methods.
\item \textbf{Genome download.} Mapped PRJNA321464 $\to$ assembly \textbf{GCA\_001931535.1} (esearch/esummary) and downloaded genome FASTA + GFF + proteome via the NCBI Datasets REST API (free, no auth). MD5 of the zip recorded.
\item \textbf{C2/C3 --- genome statistics} (\texttt{genome\_stats.py}): length + GC from FASTA; CDS/strand/tRNA/rRNA counts and coding density from GFF.
\item \textbf{C5 --- pfa cluster} (\texttt{pfa\_blast.py}): \texttt{efetch} of the 10 reference Pfa proteins the paper cites (\textit{Aetherobacter} sp.\ SBSr008 AIJ50375-77; \textit{A.\ fasciculatus} AIJ50372-74; \textit{Sorangium cellulosum} So ce56 CAN90975-77, CAN95221); \texttt{makeblastdb} + \texttt{blastp} ($e \le 10^{-10}$) vs the \textit{M.\ rosea} proteome; summed-HSP coverage per best subject; GFF synteny check of the top hits.
\item \textbf{C4 --- antiSMASH} on uicgpu: created a fresh conda env with \textbf{antiSMASH 8.0.4}, downloaded and pre-built all databases, ran \texttt{antismash --genefinding-tool prodigal --cpus 16} on the genome; parsed \texttt{mrosea.json} region areas by product category.
\item \textbf{Verdict:} two free LLM judges (Argo \texttt{gpt-5.2} and \texttt{claude-opus-4.8}) scored the claims-vs-evidence (not regex).
\end{enumerate}

\section{Results vs Paper}

\subsection{C2/C3 --- Genome statistics (Table 1) --- exact match}

\begin{longtable}{@{}l r r c@{}}
\toprule
Metric & Paper (Table 1) & This replication & Match \\
\midrule
Contigs / topology & 1, complete circular & 1 & \checkmark \\
Genome size (bp) & \textbf{16{,}040{,}666} & \textbf{16{,}040{,}666} & \checkmark\ EXACT \\
GC \% & 69.07 & 69.10 & \checkmark\ ($\Delta$0.03) \\
Protein-coding genes (CDS) & \textbf{14{,}018} & \textbf{14{,}018} & \checkmark\ EXACT \\
CDS on (+) strand & \textbf{6{,}983} & \textbf{6{,}983} & \checkmark\ EXACT \\
CDS on ($-$) strand & \textbf{7{,}035} & \textbf{7{,}035} & \checkmark\ EXACT \\
Coding density \% & 87.31 & 87.59 & \checkmark\ ($\Delta$0.28) \\
tRNA & 88 & 89 & $\sim$ ($\Delta$1) \\
rRNA operons & 4 (5S--16S--23S) & 4$\times$16S + 4$\times$23S ($+$ 2$\times$5S) & \checkmark \\
\bottomrule
\end{longtable}

The \textbf{exact reproduction of the strand-resolved CDS counts (6{,}983 / 7{,}035)} confirms this is the identical deposited assembly, and the genome-size and gene-count headline numbers reproduce to the base pair / gene. The ``largest bacterial genome'' claim is a comparative statement about 2021 databases and is not independently benchmarked here, but the 16.04\,Mbp size --- extreme for a bacterium --- is confirmed.

\subsection{C5 --- \textit{pfa} PUFA gene cluster --- reproduced (3 independent lines)}

\textbf{(a) Homology.} BLASTP of every reference Pfa protein vs the \textit{M.\ rosea} proteome (\texttt{pfa\_blast\_summary.json}):

\begin{longtable}{@{}l l r r@{}}
\toprule
Reference Pfa & Best \textit{M.\ rosea} hit & \%id & summed cov \\
\midrule
Pfa1 (\textit{Aetherobacter}/\textit{Sorangium}) & APR86155.1 & 71.3 / 64.5 & $\sim$97--99\% \\
Pfa2 / PfaA (PKS) & APR86156.1 & 67--68 & 49--79\% \\
Pfa3 / PfaC (PKS) & APR86157.1 & 56--63 & 79--91\% \\
PfaE (PPTase) & APR88149.1 & 63.1 & 98\% \\
\bottomrule
\end{longtable}

\textbf{(b) Synteny.} The three core hits are consecutive locus tags on the same strand: \texttt{A7982\_11504} (13{,}114{,}225--13{,}115{,}874 +) $\to$ \texttt{A7982\_11505} (13{,}115{,}901--13{,}123{,}181 +) $\to$ \texttt{A7982\_11506} (13{,}123{,}210--13{,}131{,}432 +), with 27 / 29 bp intergenic gaps --- a contiguous \textit{pfa} operon.

\textbf{(c) Independent annotation + antiSMASH.} NCBI PGAP independently annotates APR86156.1 as ``omega-3 polyunsaturated fatty acid synthase subunit, PfaA.'' antiSMASH independently detected a T1PKS/hglE-KS region at 13{,}095{,}900--13{,}151{,}432 (hglE-KS is the PUFA-synthase ketosynthase class) that exactly spans this operon.

\subsection{C4 --- antiSMASH BGC survey --- category-level reproduced; count version-shifted}

\begin{longtable}{@{}l r r@{}}
\toprule
Category & Paper (antiSMASH v5.0) & This replication (v8.0.4) \\
\midrule
terpene & 9 & 12 \\
RiPP-like / RiPP & 7 & 9 \\
NRPS + NRPS-like & 7 & 8 \\
PKS (T1PKS+T3PKS) & 4 & 5 \\
RRE-containing & 4 & 5 \\
indole & 3 & 3 \\
thioamitide(s) & 2 & 3 \\
arylpolyene & 2 & 2 \\
lanthipeptide & 3 & 3 \\
phosphonate & 1 & 1 \\
phenazine & 1 & 1 \\
siderophore & 1 & 1 (NI-siderophore) \\
\midrule
\textbf{TOTAL BGC regions} & \textbf{47} & \textbf{53} \\
\bottomrule
\end{longtable}

The dominant-category ranking (terpene / NRPS / RiPP / PKS) and the rare singleton set (phosphonate, phenazine, siderophore, thioamitide, arylpolyene) reproduce faithfully. The total count differs (53 vs 47); this is consistent with antiSMASH v8 having more detection rule sets than the v5 the authors used.

\section{Genuine Critique}

This section is deliberately critical --- it enumerates what did \emph{not} cleanly replicate and where the paper's inferences remain incompletely supported by the independent evidence collected here.

\subsection{What did not exactly reproduce}
\begin{itemize}[leftmargin=1.2cm]
\item \textbf{antiSMASH BGC total (C4).} The paper reports 47 BGCs; this replication finds 53. The difference is fully within the noise expected from an antiSMASH major-version jump (v5.0 $\to$ v8.0.4), which adds new detection rules (thioamitide, RRE-containing, more RiPP subclasses) and refines region boundaries. However, this replication did \emph{not} rerun antiSMASH v5.0 side-by-side to prove exact byte-for-byte parity with the paper. The claim that ``the count offset is entirely a version artifact'' is therefore plausible and consistent with the category-level agreement, but not \emph{directly} demonstrated.
\item \textbf{tRNA count.} Paper 88 vs replication 89 ($\Delta$1). Small, and typical of tRNA-scan annotation-method differences (Aragorn vs tRNAscan-SE parameter set), but not resolved.
\item \textbf{GC \% and coding density.} $\Delta$0.03 and $\Delta$0.28 respectively. Both trivial in magnitude and typical of the difference between an old RAST/manual pipeline number and a fresh Datasets/GFF calculation, but the paper does not publish enough methodological detail to reproduce the exact rounding path.
\end{itemize}

\subsection{Claims tested weakly or by proxy only}
\begin{itemize}[leftmargin=1.2cm]
\item \textbf{``Largest bacterial genome (2021)'' (C2).} This is a comparative claim about the 2021 state of NCBI/GTDB, not about \textit{M.\ rosea} itself. This replication confirms the 16.04\,Mbp size but does \emph{not} independently benchmark it against the 2021 complete-bacterial-genome distribution. It is possible other equally-large or larger deposits existed by mid-2021; this replication cannot say.
\item \textbf{Horizontal gene transfer from Actinobacteria (C5, HGT inference).} The paper argues \textit{pfa} was acquired via HGT from Actinobacteria based on synteny + phylogeny. This replication reproduces the \emph{ortholog identification and synteny}, but does \emph{not} independently redo the phylogenetic tree that grounds the HGT inference. The HGT claim is therefore consistent with the evidence gathered here but not independently corroborated at the phylogenetic-signal level.
\item \textbf{$\sim$44\% paralogous coding potential.} Mentioned in the paper's abstract; not independently recomputed in this replication (would require an all-vs-all self-BLAST + clustering step that was out of scope for this pass).
\end{itemize}

\subsection{Structural limitations of this replication}
\begin{itemize}[leftmargin=1.2cm]
\item \textbf{Single-tool antiSMASH.} Only antiSMASH v8 was run. A parallel run with PRISM, DeepBGC, or GECCO would triangulate the BGC count and category assignments; without it, the ``53 vs 47 is version-driven'' story rests on antiSMASH's own version notes.
\item \textbf{Coverage completeness of pfa BLAST.} PfaA and PfaC show summed-HSP coverage of 49--79\% and 79--91\% respectively --- strong, but not saturating. This is expected for large multi-domain PKS proteins where individual domains hit as separate HSPs, but the tabulation here does not decompose per-domain hits, so a reader cannot easily verify that \emph{every} functional domain (KS, AT, KR, DH, ACP\ldots) is present.
\item \textbf{No wet-lab confirmation.} The \textit{pfa} operon is inferred entirely from sequence + annotation + BGC-prediction convergence. No metabolite (EPA/DHA/AA) detection is performed to confirm the cluster is expressed and functional. The paper likewise infers rather than measures.
\item \textbf{LLM judges disagreed.} \texttt{gpt-5.2} voted PARTIAL (weighting the 53-vs-47 delta), \texttt{claude-opus-4.8} voted REPLICATED (attributing the delta to tool version). The verdict here follows the more conservative judge and the ``do not inflate'' guidance, but reasonable readers could land at either PARTIAL or a qualified REPLICATED.
\end{itemize}

\subsection{Net honest read}
The paper's \emph{core empirical} claims --- genome size, strand-resolved gene count, presence and synteny of the \textit{pfa} operon, dominant BGC categories --- reproduce cleanly and, in several cases, to the base pair. The paper's \emph{comparative} and \emph{inferential} claims --- ``largest bacterial genome,'' ``HGT from Actinobacteria,'' the specific 47-BGC headline count --- are either not independently benchmarked here or reproduce only at the category / structure level. That is a strong PARTIAL replication, and the verdict should not be inflated to ``full REPLICATED'' on the strength of the empirical hits alone.

\section{Evidence Artifacts}
\begin{itemize}[leftmargin=1.2cm]
\item \texttt{evidence/genome\_stats.json} --- genome statistics vs paper Table 1
\item \texttt{evidence/pfa\_blast\_summary.json} --- pfa ortholog BLAST hits + coverage
\item \texttt{evidence/antismash\_summary.json} --- antiSMASH 8.0.4 region/category tally
\item \texttt{evidence/llm\_judge\_gpt52.txt}, \texttt{evidence/llm\_judge\_opus48.txt} --- dual free-LLM judgments
\item \texttt{work/} --- fulltext.xml, downloaded genome (GCA\_001931535.1), pfa\_refs.faa, BLAST DB + outputs, antiSMASH JSON, scripts
\end{itemize}

\section{LLM-Judge Summary (free Argo endpoints)}
\begin{itemize}[leftmargin=1.2cm]
\item \textbf{gpt-5.2:} coverage 100\%, verdict \textbf{PARTIAL} (weighted the exact-count deltas, esp.\ C4 53 vs 47).
\item \textbf{claude-opus-4.8:} coverage 100\%, verdict \textbf{REPLICATED} (4 Agree + 1 Partial; C4 delta attributed to tool version).
\end{itemize}

Reconciling the two (and following the ``do not inflate'' guidance), the honest call is a \textbf{strong PARTIAL}: the genome and \textit{pfa}-cluster science reproduce robustly and largely exactly, with the single non-exact item (BGC total) explained by a documented antiSMASH version difference rather than any scientific contradiction.

\section*{Verdict}
\textbf{Verdict: PARTIAL}

\end{document}
